\documentclass[10pt,a4paper]{amsart}
\usepackage[margin=19mm]{geometry}
\usepackage{amsmath,amssymb,mathtools}
\usepackage[scr=rsfs,cal=euler]{mathalfa}
\usepackage{xfrac}
\usepackage[unicode,hidelinks,bookmarks=false]{hyperref}
\newcommand{\ie}{i.e.\ }
\newcommand{\cf}{cf.\ }
\newcommand{\resp}{respectively\ }
\newcommand{\Subsection}{\subsection}
\providecommand{\nobreakdash}{\nobreak\hbox{-}}
\setlength{\emergencystretch}{2em}
\multlinegap0pt
\usepackage{bm}
\usepackage{bbm}
\usepackage{dsfont}
\usepackage{xspace}
\usepackage{cancel}
\usepackage{mathtools}
\usepackage{amsthm}
\makeatletter
\@ifundefined{Theorem}{\newtheorem{Theorem}{Theorem}}{}
\makeatother
\makeatletter
\expandafter\ifx\csname DefConstr\endcsname\relax
\newenvironment{DefConstr}
{\pushQED{\qed}\renewcommand{\qedsymbol}{$\diamond$}\defConstr}
{\popQED\enddefConstr}
\fi
\expandafter\ifx\csname Definition\endcsname\relax
\newenvironment{Definition}
{\pushQED{\qed}\renewcommand{\qedsymbol}{$\diamond$}\definition}
\fi
\expandafter\ifx\csname Example\endcsname\relax
\newenvironment{Example}
{\pushQED{\qed}\renewcommand{\qedsymbol}{$\diamond$}\example}
\fi
\expandafter\ifx\csname Remark\endcsname\relax
\newenvironment{Remark}
{\pushQED{\qed}\renewcommand{\qedsymbol}{$\diamond$}\remark}
\fi
\newcommand{\ZZ}{{\mathbb Z}}
\makeatother
\title[The Spectral Edges Conjecture via Corners]{The Spectral Edges Conjecture via Corners}
\author{Matthew Faust, Frank Sottile}
\date{Source: arXiv:2510.10143v1 (CC BY 4.0)}
\begin{document}
\maketitle

\noindent\emph{Source and licence.} The Spectral Edges Conjecture via Corners. Version arXiv:2510.10143v1 (CC BY 4.0), distributed under the Creative Commons Attribution 4.0 International licence, as recorded in the arXiv licence metadata for this exact version and shown on the abstract page https://arxiv.org/abs/2510.10143v1. Full rights notice: licenses/arxiv-2510.10143-CC-BY-4.0.txt.\par\smallskip

\noindent\textit{Editorial scope.} Counted unit: the Theorem whose source label is \texttt{Th:FlowerGraphs} in the article (source0.tex lines 836--838, source label \texttt{Th:FlowerGraphs}), delivered together with the article's own complete original proof of it (source0.tex lines 840--892, one \texttt{proof} environment of 53 lines). No mathematics is written, repaired or re-derived: statement and proof are the article's own text line for line. Required context retained uncounted: none. Every definition, display and label of those lines is retained unchanged. Omitted from this package and not redistributed: 1--835, 839--839, 893--1469. Nothing inside a retained excerpt is omitted or altered: every retained body is delivered exactly as the article's lines are written, so no omission receipt of any kind accompanies this package. Citations: the retained excerpts carry no citation command at all, so no bibliography entry of the article is transcribed and no reference is redistributed. Reference adaptation: none was needed and none was made; every reference inside the delivered unit resolves inside it, so no unresolved reference or citation is produced. Printed slips kept literal: no printed word, symbol, hypothesis or number of the counted statement or of its proof was altered. Layout and class: the article's own class is \texttt{amsart} with packages including amsfonts, amssymb, amsmath, latexsym, mathtools, amscd, hyperref, graphicx, enumerate, xcolor, url, wrapfig, vmargin; the wrapper is the lane's pinned \texttt{amsart} editorial wrapper at 10pt on A4, which restates the theorem-like environments the retained text needs and adds the article's own macro definitions that the retained text needs (\texttt{\textbackslash ZZ}), with extra wrapper packages (bm, bbm, dsfont, xspace, cancel, mathtools). The delivered numbering is the unit's own. Assets: no retained span contains a figure environment, a graphics-inclusion command, a PDF-page inclusion, a table or a TikZ picture; no image file is copied into this package, the package declares no asset dependency and no image file is redistributed with it. Licence: version arXiv:2510.10143v1 of this article is distributed under the Creative Commons Attribution 4.0 International licence, as recorded in the arXiv licence metadata for this exact version and shown on the abstract page https://arxiv.org/abs/2510.10143v1. A full rights notice is delivered at licenses/arxiv-2510.10143-CC-BY-4.0.txt. Characters: every retained excerpt is pure ASCII, so the delivered engine, XeLaTeX with its default Latin Modern text font, reports no missing character.
\begin{Theorem}\label{Th:FlowerGraphs}
  A periodic flower graph $\Gamma_F$ is minimally sparse.
\end{Theorem}

\begin{proof}
  Observe that each vertex of $F$ represents a unique $\ZZ^d$-orbit of vertices of $\Gamma_F$, and the same for each
  edge of $F$. 
  Thus, a labeling of $\Gamma_F$ gives a labeling of $F$ and vice-versa.
  The Floquet matrix $H(z)$ of $\Gamma_F$ may be constructed directly
  from the weighted adjacency matrix $M$ of $F$.
  Indeed, the entry $H(z)_{u,u}$ of $H(z)$ is
  %
  \begin{equation}\label{Eq:uu-entry}
     V(u)\ +\ \sum_{P_i\in L} \left(z_{f(i)}+z_{f(i)}^{-1}\right)E(P_i)\,,
  \end{equation}
  %
  and if $v,w\in W$ with $(v,w)\neq (u,u)$, then 
  %
  \begin{equation}\label{Eq:H_Defn}
    H(z)_{v,w}\ \vcentcolon=\ \left\{
    \begin{array}{rcl}
      z_{f(i)}E(v,w)&&\mbox{if } v\to w\mbox{ is the directed edge in petal }P_i,\vspace{3pt}\\
      z_{f(i)}^{-1}E(v,w)&&\mbox{if } w\to v\mbox{ is the directed edge in petal }P_i,\vspace{3pt}\\
      M_{v,w}&&\mbox{otherwise.}
    \end{array}\right.
  \end{equation}

  The generic dispersion polynomial $\det(H(z){-}\lambda)$ is minimally sparse.
  Let the values of $V$ and of $E$ be indeterminates (variables).
  Note that $\det(H(z){-}\lambda)$  is obtained  by replacing each occurrence of $V(v)$ in $\det H(z)$ with $V(v){-}\lambda$.
  Thus, it suffices to show that $\det H(z)$ is minimally sparse; that the only monomials in the $z_j$ are
  single variables or their inverses.

  Writing $S_W$ for the symmetric group on the set $W$, we have
  \[
     \det H(z)\ =\ \sum_{w\in S_W} \pm \prod_{v\in W} H(z)_{v,\pi(v)}
      \ \ =\ \sum_{w\in S_W}  H(z)_\pi\,.
  \]
  The summand $H(z)_\pi$ is non-zero if and only if each entry $H(z)_{v,\pi(v)}$ for $v\in W$ is nonzero.
  When $v=\pi(v)$, so that $v$ is a fixed point of $\pi$, then $H(z)_{v,v}$ is either the constant $V(v)$ or
  else the sum~\eqref{Eq:uu-entry}, when $v=u=\pi(u)$.
  All other non-zero factors $H(z)_{v,\pi(v)}$ of $H(z)_\pi$ correspond to edges of $F$ that are not among the loops
  $L$ at $u$. 

  If  $H(z)_\pi$ is non-zero and $(v,\pi(v))$ with $v\neq \pi(v)$ is an edge in a petal $P$, then every other edge in
  $P$ occurs as $(w,\pi(w))$.
  That is, the vertices of $P$ form a cycle in $\pi$.
  This includes the distinguished vertex $u$.
  This precludes any other edge of any other petal corresponding to a factor of $H(z)_\pi$, as well as the entry $H(z)_{u,u}$.
  Similarly, if $H(z)_{u,u}$ is a factor of $H(z)_\pi$, then none of its factors correspond to edges of
  petals $P$ that are not loops in $L$. 

  These observations show that  each summand $H(z)_\pi$ has at most one factor that
  involves any variable $z_i$.
  As this factor is an entry in $H(z)$,~\eqref{Eq:uu-entry} and~\eqref{Eq:H_Defn} imply that the generic dispersion polynomial
  is minimally sparse, and completes the proof.
\end{proof}

\end{document}
